% Discussion paragraph carrying the Storyline 2 perspective inside Storylines 1 or 3.
Our results also bear on how agreement about reputations arises. Agreement is usually attributed to mechanisms that align information across observers, such as institutions, gossip and opinion synchronization\cite{radzvilavicius2021,kawakatsu2024,murase2024synchronization}, or to observers who draw on several observations at once\cite{michelmata2024}. In our model, observers share nothing and update on one observation at a time, yet they largely agree. In the mean-field theory, agreement arises because inertia makes each observer's verdict a reliable function of the target's behaviour, and independent observers who each reach a reliable verdict agree. This route complements information sharing rather than replacing it. It remains available where shared information is costly or untrusted, and it can be combined with institutions or gossip where they exist. Emergent agreement is not automatically useful, however. Under one-step inertia, observers using Shunning agree because they regard nearly everyone as bad. The direction of inertia decides whether emergent agreement is also accurate.
